\documentclass[tikz,border=8pt]{standalone}
\usepackage{ctex}
\usepackage{amsmath}
\usetikzlibrary{arrows.meta}
\begin{document}
\begin{tikzpicture}[>=Stealth]
  % ================= 左panel：几何截面 =================
  % 内柱（实心导体）
  \filldraw[fill=gray!35, draw=black, thick] (0,0) circle (0.8);
  \node at (0.25,0.2) {$\oplus$};
  \node at (-0.28,0.25) {$\oplus$};
  \node at (0.0,-0.32) {$\oplus$};
  \node[below=2pt] at (0,-0.82) {内柱 $+\lambda$};
  % 外筒（薄导体筒）
  \draw[very thick] (0,0) circle (2.2);
  \node at (90:2.2) {$\ominus$};
  \node at (150:2.2) {$\ominus$};
  \node at (210:2.2) {$\ominus$};
  \node at (270:2.2) {$\ominus$};
  \node at (330:2.2) {$\ominus$};
  \node[above=2pt] at (0,2.22) {外筒 $-\lambda$};
  % 间隙电场箭头（沿径向向外）
  \foreach \ang in {60,120,185,235,300}
    \draw[->,teal,thick] (\ang:0.95) -- (\ang:2.02);
  \node[teal] at (30:1.5) {$\vec{E}$};
  % 水平虚线半径 + 三个区域标记
  \draw[dashed,gray] (0,0) -- (2.75,0);
  \fill (0.4,0) circle (1.6pt);
  \node[below] at (0.4,-0.06) {\footnotesize\textcircled{\scriptsize 1}};
  \fill (1.5,0) circle (1.6pt);
  \node[below] at (1.5,-0.06) {\footnotesize\textcircled{\scriptsize 2}};
  \fill (2.75,0) circle (1.6pt);
  \node[below] at (2.75,-0.06) {\footnotesize\textcircled{\scriptsize 3}};
  \node[below=8pt] at (0.8,0) {$a$};
  \node[above right=2pt] at (2.2,0.02) {$b$};
  % 中心点
  \fill (0,0) circle (1.6pt);
  % 参考点注记（与球形对的关键差异；置于底部避免与外筒相碰）
  \node[font=\footnotesize,align=left,anchor=west] at (-3.4,-2.95)
    {参考点：$V=0$ 取在 $r=b$（线电荷电势在无穷远发散）};

  \begin{scope}[xshift=8.2cm]
    % ================= 右panel：V(r) 剖面 =================
    \draw[->] (0,0) -- (3.45,0) node[below] {$r$};
    \draw[->] (0,0) -- (0,1.45) node[left] {$V(r)$};
    % 平台段 ①：r<a，平台高度 = ln(b/a) 归一化为 1
    \draw[very thick,blue] (0.12,1) -- (1,1);
    % 对数下降段 ②：a<r<b, V/V_ab = ln(2.75/x)/ln(2.75)
    \draw[very thick,blue,domain=1:2.75,samples=80,smooth]
      plot (\x,{ln(2.75/\x)/1.012});
    % 零段 ③：r>b
    \draw[very thick,blue] (2.75,0) -- (3.3,0);
    % 参考虚线
    \draw[dashed,gray] (0,1) -- (1,1);
    \draw[dashed,gray] (1,0) -- (1,1);
    \node[left] at (0,1) {$V_{ab}$};
    \node[below] at (1,-0.02) {$a$};
    \node[below] at (2.75,-0.02) {$b$};
    % 区域标记
    \node[font=\footnotesize] at (0.55,1.14) {\textcircled{\scriptsize 1}};
    \node[font=\footnotesize] at (1.75,0.55) {\textcircled{\scriptsize 2}};
    \node[font=\footnotesize,above] at (3.0,0.04) {\textcircled{\scriptsize 3} $V=0$};
    % 公式注记（置于曲线上方空白处）
    \node[font=\footnotesize,anchor=west] at (1.7,0.72)
      {$\frac{\lambda}{2\pi\varepsilon_0}\ln\frac{b}{r}$};
  \end{scope}
\end{tikzpicture}
\end{document}
